\documentclass[hyphens,aspectratio=169,dvipsnames]{beamer}
\usepackage{graphicx}
\usepackage{xcolor}
\usepackage{amssymb}
\usepackage{amsmath}
\usepackage{mathtools}
\usepackage{textcomp}
\usepackage{moresize}
\usepackage{framed}
\usepackage{minted}
\usepackage{relsize}
\usepackage{emoji}
\usepackage{tikz}
\usepackage{comment}
\usetikzlibrary{shapes.geometric, arrows, positioning}
\usetheme{Berlin}
\usecolortheme[RGB={0,95,47}]{structure}
\beamertemplatenavigationsymbolsempty

\makeatletter
\AtBeginEnvironment{minted}{\dontdofcolorbox}
\def\dontdofcolorbox{\renewcommand\fcolorbox[4][]{##4}}
\makeatother

\newcommand{\textpf}[1]{\texttt{\color{black}\fcolorbox{lightgray}{lightgray}{#1}}}

\begin{document}

\begin{frame}[fragile]{Topics}
    \begin{itemize}
        \item Threads
    \end{itemize}
\end{frame}

\begin{frame}[fragile]{\textpf{fork}ing is a heavyweight operation}
    \begin{itemize}
        \pause \item \textpf{fork} duplicates (almost) the entire state of the process!
        \pause \item Then, because processes don't share an address space, you need to use IPC to communicate between them.
        \pause \item Sometimes, this is overkill. What if I just want to do multiple things at once \textit{within a process}?
    \end{itemize}
\end{frame}

\begin{frame}[fragile]{Threads}
    \begin{itemize}
        \pause \item Threads are execution units \textit{within a process}.
        \pause \item That is, a process with multiple threads can do multiple things at the same time.
        \pause \item Threads exist \textit{within a process}, so they share an address space.
    \end{itemize}
\end{frame}

\begin{frame}[fragile]{Creating a Thread}
    \begin{itemize}
        \pause \item When a process starts, it has only one thread (the \textit{main thread}).
        \pause \item To create a new thread, use the \textpf{thrd\_create} function from \textpf{threads.h}.
        \pause \item (Run \textpf{info libc 'ISO C Threads'} to read the documentation for the functions in \textpf{threads.h}.)
    \end{itemize}
\end{frame}

\begin{frame}[fragile]{New Terminology: ``Handle"}
    \begin{itemize}
        \pause \item A handle is an object that represents a reference to a resource.
        \pause \item Handles can be passed to code that operate on the underlying resource.
        \pause \item For example,
        \begin{itemize}
            \pause \item A file descriptor is a handle for a socket, I/O stream, or file. % It's more complicated than this, but it's okay :)
            \pause \item A PID is a handle for a process.
        \end{itemize}
        \pause \item The handle for a thread is called a \textpf{thrd\_t}.
    \end{itemize}
\end{frame}

\begin{frame}[fragile]{\textpf{thrd\_create}}
    \begin{itemize}
        \pause \item \textpf{thrd\_create} takes three arguments:
        \begin{enumerate}
            \pause \item \textpf{thrd\_t *thr}, an out-argument into which your thread handle is written.
            \pause \item \textpf{thrd\_start\_t func}, the function that your thread will execute. (Should take a \textpf{void *} and return an \textpf{int}).
            \pause \item \textpf{void *arg}, the argument to \textpf{func}.
        \end{enumerate}
        \pause \item Q: Why is \textpf{arg} a \textpf{void *}?
        \pause \item On success, \textpf{thrd\_create} returns \textpf{thrd\_success}.
    \end{itemize}
\end{frame}

\begin{frame}[fragile]{Demo}
    Demo :)
\end{frame}

\begin{frame}[fragile]{Mutexes}
    \begin{itemize}
        \pause \item Sometimes, there's a \textit{critical section} of your program that should be executed by only one thread at a time.
        \pause \item You can synchronize threads to enforce this constraint using a \textit{mutex}.
        \pause \item A mutex is an object with two states:
            \begin{itemize}
                \pause \item \textit{locked}, indicating that the mutex is owned by a thread.
                \pause \item \textit{unlocked}, indicating that it's not.
            \end{itemize}
        \pause \item When two threads simultaneously attempt to lock a mutex, only one of them will get it.
        \pause \item The other one will be blocked until the mutex is unlocked.
    \end{itemize}
\end{frame}

\begin{frame}[fragile]{Mutex Functions}
    \begin{itemize}
        \pause \item \textpf{mtx\_init}: initializes a mutex. Must be called before a mutex can be used.
        \pause \item \textpf{mtx\_lock}: locks the mutex.
        \pause \item \textpf{mtx\_unlock}: unlocks the mutex.
    \end{itemize}
\end{frame}

\begin{frame}[fragile]{Demo, continued!!}
    Demo :)
\end{frame}

\begin{frame}[fragile]{\textpf{thrd\_join}}
    \begin{itemize}
        \pause \item The equivalent of \textpf{waitpid} for threads is called \textpf{thrd\_join}.
        \pause \item \textpf{thrd\_join} takes two arguments:
            \begin{enumerate}
                \pause \item the \textpf{thrd\_t} to be waited upon
                \pause \item an out-argument \textpf{int *} to store returned value from the thread's start function.
            \end{enumerate}
    \end{itemize}
\end{frame}


\begin{frame}[fragile]{Activity!}
    \begin{itemize}
        \item Write a program that uses \textpf{getrandom(2)} to read a gigabyte of random data into a static array, and return its sum (mod 256).
        \item Use threads to divide and conquer the work.
        \item Measure the performance benefit of the threaded program using \textpf{time(1)}.
    \end{itemize}
\end{frame}

\end{document}
